\documentclass[11pt]{article}
\usepackage[a4paper,margin=1in]{geometry}
\usepackage[T1]{fontenc}
\usepackage{lmodern,microtype}
\usepackage{amsmath,amssymb,amsthm,mathtools}
\usepackage{booktabs,enumitem}
\usepackage{xcolor}
\usepackage[numbers,sort&compress]{natbib}
\usepackage[colorlinks=true,linkcolor=blue!55!black,citecolor=blue!55!black]{hyperref}

\newtheorem{theorem}{Theorem}[section]
\newtheorem{proposition}[theorem]{Proposition}
\newtheorem{corollary}[theorem]{Corollary}
\newtheorem{lemma}[theorem]{Lemma}
\theoremstyle{remark}
\newtheorem{remark}[theorem]{Remark}

\title{Extended Deterministic Numerator Anchor Atlas}
\author{Bin Wang}
\date{July 2026}

\begin{document}
\maketitle

\begin{abstract}
The analytic-tail theorem of RH-252 is exact but does not identify the
finite target coefficients beyond the order-12 table of RH-243.  We evaluate
the same deterministic periodic-point formula through order $28$.  This adds
16 orders and enumerates $32767$ physical fixed points at the largest order.
The new order block contributes $0.0021942543$ to the unit-disk logarithmic
norm, and its descriptive log-linear root rate is $0.7009986$.  These are
finite reproducible facts.  They do not prove a uniform trace envelope, a
cloud coefficient bridge, or any statement about Riemann zeros.
\end{abstract}

\section{Exact target dictionary}

Let $P_n$ denote the deterministic flat periodic trace from the quadratic
band-merging map.  The RH-243 dictionary is
\begin{equation}
 a_n=r_H^{-n}\left[P_n-1-(-1)^n+2\mathbf 1_{2\mid n}\lambda^{-n/2}\right],
 \qquad r_H=0.85,
 \label{eq:anchor}
\end{equation}
with $\lambda=1.6785735104283177\ldots$ \citep{WangRH243}.  The periodic
point and multiplier formulas used to evaluate $P_n$ are the exact finite
formulas archived with the Collet--Eckmann completion \citep{WangRH11}.

\begin{proposition}[Finite extension]
\label{prop:finite}
For every integer $2\le n\le28$, equation \eqref{eq:anchor} defines an exact
finite deterministic target value.  The order-$28$ computation enumerates
$32767$ distinct physical fixed points before summing their flat weights.
\end{proposition}

\begin{proof}
The inverse-branch construction in RH-11 enumerates all cyclic Markov words,
removes the duplicated shared endpoint, and verifies the physical fixed-point
count.  Applying the displayed finite sum for each order and then the explicit
algebraic correction in \eqref{eq:anchor} gives the claim.
\end{proof}

\section{The new order block}

The following values are representative rows from the 16 new orders.  The
last column is the Hardy-scaled target, not a noisy trace.

\begin{center}
\small
\begin{tabular}{rrrr}
\toprule
$n$ & $P_n$ & unscaled $g_n$ & $a_n$\\
\midrule
13 & 0.00118921 & 0.00118921 & 0.00983586\\
16 & 1.96851809 & 0.000250789 & 0.00337758\\
20 & 1.98876937 & 0.0000316550 & 0.000816703\\
24 & 1.99600689 & 0.00000399130 & 0.000197269\\
28 & 1.99858189 & 0.000000502810 & 0.0000476073\\
\bottomrule
\end{tabular}
\end{center}

For the unit disk and the logarithmic norm
\begin{equation}
 J_{I}(1):=\sum_{n\in I}\frac{|a_n|}{n},
 \label{eq:norm}
\end{equation}
the archived computation gives
\begin{align}
 J_{2:12}(1)&=0.49450543569144195,\\
 J_{13:28}(1)&=0.0021942543215719553,\\
 J_{2:28}(1)&=0.496699690013014.
 \label{eq:values}
\end{align}
At radius $0.8$, the new block contributes only
$8.7150803572\times10^{-5}$.

\section{Finite slope diagnostic and its boundary}

Fitting $\log|a_n|$ linearly against $n$ on the new block gives
\begin{equation}
 \widehat q_{13:28}=0.7009986349256669,
 \quad
 \widehat q_{\rm odd}=0.7009206511866494,
 \quad
 \widehat q_{\rm even}=0.7011096165507584.
 \label{eq:fit}
\end{equation}
The agreement between parity subsequences is a useful finite consistency
check with the RH-252 radius budget \citep{WangRH252}.  It is not a proof of
$|a_n|\le Cq^n$ for all $n$: a finite fit has no control over orders $29$ and
beyond, nor over the constant needed in the Cauchy tail theorem.

\begin{remark}[Scope]
The atlas is deterministic and finite.  It does not select a noisy spectral
cloud, certify a quotient block, or identify the coefficients of a moving
operator.  The target is an anchor, not an observed eigenvalue list.
\end{remark}

\section{Route decision}

RH-252 provides an all-order analytic existence statement for the target tail,
while this paper supplies a longer finite target table.  The next genuinely
new input must come from the spectral side: an expanded resolved candidate
window or an invariant quotient selector.  Reusing the frozen RH-248 shell
window would not be progress.  Gates A--E remain false/open; no
Hilbert--Polya operator, zeta-divisor equality, Riemann-zero identification,
or RH implication is asserted.

\bibliographystyle{plainnat}
\bibliography{references}
\end{document}
